\section{Theoretical Foundation and Literature Review}
\subsection{Shortest-Path Search and Heuristic Assumptions}
Once a graph is available, both applications require minimum-cost paths. On directed $G=(V,E,w)$ with nonnegative weights, Dijkstra orders pending states by tentative start cost~\cite{dijkstra1959}. \astar{} adds estimated remaining cost to guide search toward the goal~\cite{hart1968}:
\begin{equation}f(n)=g(n)+h(n).
\end{equation}
Here $g(n)$ is the best start-to-$n$ cost found so far and $h(n)$ estimates cost to goal $t$. OPEN contains pending states; improvements update costs and predecessors. The goal is accepted on valid minimum-priority extraction, not discovery. Setting $h=0$ gives Dijkstra in the same implementation.

This stopping rule requires suitable estimates. Admissibility means $0\leq h(n)\leq h^*(n)$, where $h^*(n)$ is minimum remaining graph cost. Consistency additionally requires
\begin{equation}
h(u)\leq w(u,v)+h(v)\quad\text{for every }(u,v)\in E,\qquad h(t)=0.
\label{eq:consistency}
\end{equation}
Consistency makes the first valid expansion's cost optimal. Admissible but inconsistent estimates can require reopening improved states. These distinctions and tie resolution underlie Hart et al.~\cite{hart1968} and the implementation in Section~5. Fewer expansions need not mean lower wall time: heuristic evaluation and representation also incur work.

\subsection{Geometric Representation and Replanning}
Optimality depends on admitted paths. Grids restrict headings and may miss narrow passages; visibility-based planning uses obstacle-free geometric connections~\cite{lozano1979}. Finite footprints require configuration-space treatment. A complete polygonal visibility graph and a sparse scan-center graph admit different paths, so optimality in the latter need not imply a continuous shortest path.

Theta* uses grid search with off-grid connections~\cite{daniel2010}. Experiments on random grids and game maps compare length, time, expansions, and heading changes, finding shorter paths than grid smoothing without a full visibility graph. Basic Theta* nevertheless does not guarantee a continuous shortest path. Our ordinary A* on verified continuous links likewise requires separate assessment of representation and search.

When observations change, D* Lite reuses search information through locally consistent value updates~\cite{koenig2002}. Its unknown-terrain navigation and mapping experiments reduce operations relative to repeated search. Those grids initially allow unknown traversal optimistically; our graph admits only verified free links. Repeated A* preserves implementation transparency. Incremental replanning remains an alternative whose project performance is unmeasured.

\subsection{Observation, Exploration, and Escape}
Constructing links first requires representing sensor knowledge. Elfes accumulates uncertain evidence in probabilistic occupancy grids~\cite{elfes1989}; sonar and stereo demonstrations show the value of additional viewpoints. Our deterministic labels and polygon unions accumulate observations and preserve unknown space, but omit probabilistic uncertainty and localization drift.

Free/unknown boundaries then suggest exploration targets. Yamauchi visits accessible frontiers~\cite{yamauchi1997}; office-robot experiments address sensor artifacts, while idealized coverage remains conditional on sensing and reachability. Our polygon candidates similarly seek observations, using ancestry for recovery and C++ A* for routing. Estimated utility does not prove information gain, and one unreachable candidate does not exhaust a branch.

TangentBug combines range discontinuities and a local tangent graph with target motion and boundary following~\cite{kamon1996}. Its convergence analysis assumes static finite-perimeter geometry and ideal point motion; selected simulations compare it with VisBug. Our simpler observed-parent policy does not inherit that guarantee. Together these studies explain why local path optimality requires a separate exploration and escape mechanism.

\subsection{Bundle-Based Planning and the Assigned BAR Problem}
Phan et al.~\cite{phan2026} accumulate sights, rank active open points, construct a skeleton graph, and shorten polylines through ordered bundles of line segments. DAP adjusts vertices on critical segments within represented explored space; Section 3 relates approximation quality to geometric and numerical assumptions. Section 4.5 constructs bundles and Algorithm 2 integrates observations, ranking, refinement, and movement. The main contribution is geometric representation and refinement, not ordinary graph search.

Algorithm 2 is explicitly not generally convergent. Section 5 studies a special algorithm with obstacle-diameter and near-goal connectivity assumptions, which do not prove completion of our maze exploration. Section 6 compares local paths with a grid A* baseline and RRT*, and further evaluates RRTX. Dijkstra is discussed as a foundational search method rather than evaluated as an independent experimental baseline. The source's \source{Graph.BFS_skeleton_path} performs Dijkstra-style cumulative-cost priority-queue search, despite its BFS-oriented name; this skeleton search precedes bundle-based refinement.

These differences determine the scope of our comparison: observations and endpoints are held fixed, while graph representations, boundary rules, and motion semantics are disclosed. The application addresses exploration with limited sensing and conditionally certified return to observed ancestors. It reproduces neither bundle optimization nor general polygonal BAR detection. The supplementary literature matrix records the relevance of all eight supplied PDFs, representing seven unique studies.
